\documentclass[12pt, a4paper]{article}
\usepackage[utf8]{inputenc}
\usepackage[spanish, es-tabla]{babel} % es-tabla para que diga 'Tabla' en vez de 'Cuadro'
\usepackage{amsmath}
\usepackage{graphicx}
\usepackage{geometry}
\usepackage[american]{circuitikz} % Especificamos estilo americano globalmente
\usetikzlibrary{babel} % CORRECCIÓN IMPRESCINDIBLE: Evita conflictos con flechas -> y babel español
\usepackage{float}
\usepackage{array}
\usepackage{enumitem} % Para mejor control de listas
\usepackage{comment}    

% Configuración de márgenes
\geometry{a4paper, margin=1in}

% Configuración de hipervínculos (Debe ir al final del preámbulo usualmente)
\usepackage[colorlinks=true, linkcolor=blue, urlcolor=blue, citecolor=black, filecolor=black]{hyperref}

% Redefinición de nombres (aunque babel con es-tabla hace gran parte)
\addto\captionsspanish{%
  \renewcommand{\contentsname}{Índice}%
  \renewcommand{\listfigurename}{Índice de Figuras}%
  \renewcommand{\listtablename}{Índice de Tablas}%
}

% --- Profundidad del Índice ---
\setcounter{tocdepth}{3} % 1=section, 2=subsection, 3=subsubsection

\begin{document}

% Portada o Título (Opcional, se puede agregar aquí)

% --- ÍNDICES ---
{
  \hypersetup{linkcolor=black} % Para que el índice no se vea todo azul, solo funcional
  \tableofcontents
  \newpage
  \listoffigures
  \newpage
  \listoftables
  \newpage
}

% ================= PRÁCTICA 1 =================
\phantomsection % Anclaje para el hipervínculo
\section*{\centering PRÁCTICA 1 \\ MANEJO DEL INSTRUMENTAL ELÉCTRICO Y DETERMINACIÓN DE LA 
CARACTERÍSTICA CORRIENTE-VOLTAJE (I-V) DE ALGUNOS DISPOSITIVOS}
\addcontentsline{toc}{section}{PRÁCTICA 1: MANEJO DEL INSTRUMENTAL ELÉCTRICO...}
 
\phantomsection
\subsection*{1. OBJETIVOS}
\addcontentsline{toc}{subsection}{1. OBJETIVOS}
\begin{itemize}
    \item Lograr que el estudiante aprenda a utilizar los instrumentos eléctricos de medida y a 
    realizar montaje de circuitos sencillos.
    \item Obtener la característica I-V de algunos dispositivos y determinar si cumplen o no la ley 
    de Ohm.
\end{itemize}

\phantomsection
\subsection*{2. PARTE TEÓRICA}
\addcontentsline{toc}{subsection}{2. PARTE TEÓRICA}
Esta teoría es común para las prácticas 'Determinación de una resistencia eléctrica' y 
'Determinación de la resistencia interna de un Voltímetro, un Amperímetro y una Fuente'.

\subsubsection*{2.1. RESISTENCIA, RESISTIVIDAD Y CONDUCTIVIDAD}
\addcontentsline{toc}{subsubsection}{2.1. RESISTENCIA, RESISTIVIDAD Y CONDUCTIVIDAD}
La resistencia eléctrica es la propiedad de un material que caracteriza su oposición al paso de 
una corriente eléctrica. En términos de la diferencia de potencial V, aplicada a sus extremos y de 
la intensidad de corriente I, que circula a través del conductor, la resistencia R, la definimos 
así:
\begin{equation}
    R = \frac{V}{I}
    \label{eq:ohm_basica}
\end{equation}
La unidad en el S.I. de resistencia es el ohm, que se representa por la letra griega $\Omega$. 
Según [1], $1\Omega = \frac{1V}{4}$.

Una cantidad relacionada con R es la resistividad $\rho$, o resistencia específica, la cual es 
la característica de un material y no de una muestra específica del mismo. Para materiales 
isótropos, la resistividad se define como:
\begin{equation}
    \rho = \frac{E}{J}
    \label{eq:resistividad}
\end{equation}
donde E es el campo eléctrico y J es la densidad de corriente o corriente por unidad de área de 
la sección transversal del conductor es decir:
\begin{equation}
    J = \frac{I}{A}
    \label{eq:densidad_corriente}
\end{equation}
Con el fin de expresar la resistencia R, en términos de la naturaleza del material (resistividad) y 
de la geometría del conductor, consideremos un conductor cilíndrico de longitud $l$ y sección 
transversal de área $A$ por el cual circula una corriente I al aplicarse entre sus extremos una 
diferencia de potencial V, tal como se muestra en la figura~\ref{fig:conductor}.

Si suponemos el campo eléctrico constante debe cumplirse que:
\begin{equation}
    E = \frac{V}{l}
    \label{eq:campo_electrico}
\end{equation}
Sustituyendo los valores de las ecuaciones~\eqref{eq:densidad_corriente} 
y~\eqref{eq:campo_electrico} en la ecuación~\eqref{eq:resistividad}, obtenemos que:
\begin{equation}
    \rho = \frac{RA}{l}
    \label{eq:resistividad_geo}
\end{equation}
Esta relación nos lleva a una expresión para R dada por:
\begin{equation}
    R = \rho \frac{l}{A}
\end{equation}
Según la cual R depende de la naturaleza y de la geometría del conductor.

\begin{figure}[H]
    \centering
    % Placeholder visual usando TikZ para simular el conductor
    \begin{tikzpicture}
        \draw[left color=gray!20, right color=gray!20, middle color=gray!50] (0,0) ellipse (0.2 and 0.5);
        \draw[left color=gray!20, right color=gray!20, middle color=gray!50] (4,0) ellipse (0.2 and 0.5);
        \draw[top color=gray!20, bottom color=gray!50, middle color=gray!30] (0,0.5) -- (4,0.5) arc (90:-90:0.2 and 0.5) 
            -- (0,-0.5) arc (-90:90:0.2 and 0.5);
        \draw[->] (-0.5,0) -- (0.5,0) node[above] {$I$};
        \draw[<->] (0,-0.7) -- (4,-0.7) node[midway, below] {$l$};
        \node at (4.5, 0) {$A$};
    \end{tikzpicture}
    \caption{Tramo de conductor de longitud l y sección transversal de área A, a través del cual 
    circula una corriente I.}\label{fig:conductor}
\end{figure}

De \eqref{eq:resistividad} se define la unidad en el S.I. de resistividad como: $\Omega \cdot m$.

La conductividad es un parámetro relacionado con el inverso de la resistividad. Generalmente se denota con la letra griega 
$\sigma$ y se define como:

\begin{equation}
    \sigma = \frac{1}{\rho} = \frac{J}{E}
    \label{eq:conductividad}
\end{equation}
En término de esta cantidad, la densidad de corriente se define también como:
\begin{equation}
    \vec{J} = \sigma \vec{E}
    \label{eq:ley_ohm_vectorial}
\end{equation}
De la ecuación~\eqref{eq:ley_ohm_vectorial} se deduce que J es un vector definido como el producto de un escalar $\sigma$ 
por el vector intensidad de campo eléctrico ($\vec{E}$), esto es:
\begin{equation}
    J = \sigma E
\end{equation}

\subsubsection*{2.2. LEY DE OHM}
\addcontentsline{toc}{subsubsection}{2.2. LEY DE OHM}
En 1820 George Simon Ohm (1787-1854), demostró experimentalmente que la resistencia R de ciertos conductores (los metales) 
bajo condiciones de temperatura constante, es independiente de la diferencia de potencial V aplicada en sus extremos y de 
la intensidad de corriente I que circula a través de ellos. En forma de ecuación, lo anterior se expresa así:
\begin{equation}
    R = \frac{V}{I} = \text{constante}
\end{equation}
Debido a que esta ley sólo se cumple para los metales, se dice que la misma no describe una propiedad general de la materia, 
sino una característica particular de ciertas sustancias.

Los conductores que cumplen la ley de Ohm se denominan óhmicos o lineales (su característica I-V es una recta), mientras que 
los que no la obedecen se denominan no lineales. Como ejemplo de estos últimos tenemos los materiales semiconductores.

\subsubsection*{2.3. RESISTORES}
\addcontentsline{toc}{subsubsection}{2.3. RESISTORES}
Un resistor es un componente de cualquier circuito, cuyo propósito es el de tener resistencia. Los materiales usados como 
resistores se caracterizan por tener un coeficiente de temperatura (CT) pequeño; o sea, la variación porcentual de la 
resistencia por grado de temperatura es pequeña. Los valores de la resistencia, del coeficiente de temperatura y la máxima 
potencia permitida en un resistor de algunos de los conductores más comunes, se muestran en la tabla 1.

\begin{table}[H]
    \centering
    \small
    \caption{Características de resistores comunes.}
    \begin{tabular}{|l|c|c|c|}
        \hline
        \textbf{TIPO} & \textbf{RANGO} & \textbf{C. T. (\%/°C)} & \textbf{POTENCIA MÁX} \\
        \hline
        CARBÓN & $1\Omega - 22M\Omega$ & 0.1 & 2W \\
        \hline
        ALAMBRE & $1\Omega - 100K\Omega$ & 0.0005 & 200W \\
        \hline
        PELÍCULA METÁLICA & $0.1\Omega - 10^{12}\Omega$ & 0.0001 & 1W \\
        \hline
        PELÍCULA DE CARBÓN & $10\Omega - 100M\Omega$ & Entre 0.015 - 0.05 & 2W \\
        \hline
    \end{tabular}
\end{table}

En el caso de los resistores de carbón, existe un código de colores y una clave para determinar el valor de la resistencia. 
La clave consiste en leer los anillos o bandas de colores de izquierda a derecha (el extremo es el más próximo a la primera banda) 
de tal manera que las dos primeras bandas, AB, definen las dos primeras cifras del valor de la resistencia; la tercera banda, 
C, es un multiplicador y la cuarta banda, D, indica la tolerancia o error en la determinación de R. Así pues, el valor de la 
resistencia del resistor se especifica como:
\[ R = AB \cdot 10^C \pm D \]
Para determinar el valor del resistor debemos identificar los colores según lo indica la tabla 2.

\begin{table}[H]
    \centering
    \caption{Código de colores para resistores.}
    \begin{tabular}{|l|c|c|c|}
        \hline
        \textbf{COLOR} & \textbf{VALOR} & \textbf{POTENCIA DE 10} & \textbf{TOLERANCIA} \\
        \hline
        NEGRO & 0 & $10^0$ & $\pm 20\%$ \\
        \hline
        MARRÓN & 1 & $10^1$ & $\pm 1\%$ \\
        \hline
        ROJO & 2 & $10^2$ & $\pm 2\%$ \\
        \hline
        ANARANJADO & 3 & $10^3$ & $\pm 3\%$ \\
        \hline
        AMARILLO & 4 & $10^4$ & --- \\
        \hline
        VERDE & 5 & $10^5$ & $\pm 5\%$ \\
        \hline
        AZUL & 6 & $10^6$ & --- \\
        \hline
        VIOLETA & 7 & $10^7$ & --- \\
        \hline
        GRIS & 8 & $10^8$ & --- \\
        \hline
        BLANCO & 9 & $10^9$ & --- \\
        \hline
        PLATA & --- & $10^{-2}$ & $\pm 10\%$ \\
        \hline
        ORO & --- & $10^{-1}$ & $\pm 5\%$ \\
        \hline
        NINGUNO & --- & --- & $\pm 20\%$ \\
        \hline
    \end{tabular}
\end{table}

A manera de ejemplo, supongamos que nos dan dos resistores con la siguiente secuencia de colores:
\begin{itemize}
    \item[a.] Amarillo, Violeta, Marrón, Plata: $47 \times 10^1 = 470\Omega \pm 10\%$
    \item[b.] Marrón, Negro, Rojo, Ninguno: $10 \times 10^2 = 1k\Omega \pm 20\%$
\end{itemize}
De acuerdo con el código de colores de la tabla 2, los resistores tendrían valores:
\begin{itemize}
    \item[a.] $(470 \pm 10\%) \Omega$.
    \item[b.] $(1000 \pm 20\%) \Omega$.
\end{itemize}

\subsubsection*{2.4. POTENCIA MÁXIMA DE UN RESISTOR}
\addcontentsline{toc}{subsubsection}{2.4. POTENCIA MÁXIMA DE UN RESISTOR}
Siempre que una corriente eléctrica de intensidad i circula por un resistor R, se disipa energía en forma de calor 
(efecto Joule) en el resistor. La tasa de producción de calor o potencia disipada es $P=RI^2$. Esto hace que suba 
la temperatura en el resistor y si no se evita sobrepasar el límite de temperatura máxima, se puede destruir el 
resistor. Por lo tanto, es muy importante tener presente la potencia máxima, $P_{\text{máx.}}$, que pueden soportar  
los resistores de carbón cuando se le va a conectar en un circuito. A partir de este valor máximo de P, se puede 
determinar la máxima corriente $I_{\text{máx.}}$, o la máxima caída de potencial $V_{\text{máx.}}$ que puede soportar 
el resistor sin dañarse empleando la relación:

\begin{equation}
    P_{\text{máx.}} = RI_{\text{máx.}}^2 = \frac{V_{\text{máx.}}^2}{R}
\end{equation}

\subsubsection*{2.5. AMPERÍMETROS, VOLTÍMETROS Y OHMÍMETROS}
\addcontentsline{toc}{subsubsection}{2.5. AMPERÍMETROS, VOLTÍMETROS Y OHMÍMETROS}
Las cantidades V, R, e I son cantidades macroscópicas de interés fundamental cuando se trabaja 
efectuando medidas eléctricas en conductores reales. Por lo tanto, son las cantidades que 
normalmente se miden con instrumentos debidamente calibrados. Estos instrumentos son: el 
amperímetro, para medir la intensidad de corriente I; el voltímetro, para medir la diferencia de 
potencia V y el ohmímetro, para medir la resistencia R.

\phantomsection
\subsection*{3. PARTE EXPERIMENTAL}
\addcontentsline{toc}{subsection}{3. PARTE EXPERIMENTAL}

\subsubsection*{3.1 MATERIAL}
\addcontentsline{toc}{subsubsection}{3.1 MATERIAL}
\begin{itemize}
    \item 1 fuente de poder de 0 – 30 v, 6 A
    \item 2 multimedidores Keithley (179, 175 ó 175-A)
    \item 1 diodo semiconductor
    \item 1 Resistencia (Variable o de carbón)
\end{itemize}

\subsubsection*{3.2. PROCEDIMIENTO}
\addcontentsline{toc}{subsubsection}{3.2. PROCEDIMIENTO}
Lea detalladamente el material de instrucción de los multimedidores digitales Keithley 179 y 175 en cada uno de los rangos como 
amperímetro, voltímetro (para medida en DC y AC) y Ohmímetro.

\paragraph{3.2.1 CARACTERÍSTICA I-V DE UN RESISTOR.}
\addcontentsline{toc}{subsubsection}{3.2.1 CARACTERÍSTICA I-V DE UN RESISTOR.}
\begin{itemize}
    \item Realice el montaje del circuito mostrado en la figura~\ref{fig:circuito_resistor}.
    \item Variando la salida de la fuente, para un valor fijo en el resistor R, obtenga una serie de 
        valores (20 mínimos) de V e I.
\end{itemize}

\begin{figure}[H]
    \centering
    \begin{circuitikz}[american]
        % Rama principal: Fuente -> Amperímetro -> Resistor
        \draw (0,0) node[ground]{} 
              to[V, l=$E$, invert] (0,3)
              to[ammeter, l=$A$] (3,3) -- (3,3) % Amperímetro en serie
              to[R, l=$R$, v=$V_R$] (3,0) -- (0,0); % Resistor
        
        % Rama del Voltímetro (en paralelo al Resistor)
        \draw (3,3) -- (5,3)
              to[voltmeter, l=$V$] (5,0)
              -- (3,0);

        \node[ocirc] at (0,0) {};
        \node[ocirc] at (3,0) {};
        \node[ocirc] at (3,3) {};
    \end{circuitikz}
    \caption{Circuito para el Experimento 3.2.1}\label{fig:circuito_resistor}
\end{figure}

\paragraph{3.2.2 CARACTERÍSTICA I-V DE UN DIODO SEMICONDUCTOR.}
\addcontentsline{toc}{subsubsection}{3.2.2 CARACTERÍSTICA I-V DE UN DIODO SEMICONDUCTOR.}
\begin{itemize}
    \item Realice el montaje del circuito mostrado en la figura~\ref{fig:circuito_diodo}.
    \item Variando la salida de la fuente o variando el resistor R, obtenga una serie de valores 
    (20 mínimos) de V e I.
\end{itemize}

\begin{figure}[H]
    \centering
    \begin{circuitikz}[american]
        % Rama Izquierda: Fuente
        \draw (0,0) node[ground]{} 
              to[V, l=$E$, invert] (0,3); 
        
        % Rama Superior: Resistencia Variable en serie con Amperímetro
        \draw (0,3) to[potentiometer, l=$R$] (3,3)
                    to[ammeter, l=$A$] (5,3);

        % Rama Derecha: Diodo
        \draw (5,3) to[D, l=$D$] (5,0) 
                    -- (0,0); 
        
        % Rama Externa: Voltímetro en paralelo con el Diodo
        \draw (5,3) -- (7,3) 
                    to[voltmeter, l=$V$] (7,0) 
                    -- (5,0);

        \node[ocirc] at (0,0) {};
        \node[ocirc] at (5,3) {};
        \node[ocirc] at (5,0) {};
    \end{circuitikz}
    \caption{Circuito con Resistencia Variable en Serie}
    \label{fig:circuito_diodo}
\end{figure}

\phantomsection
\subsection*{4. ACTIVIDADES}
\addcontentsline{toc}{subsection}{4. ACTIVIDADES}
\begin{itemize}
    \item Para cada experimento construya la gráfica I Vs V, y escriba la conclusión respectiva.
\end{itemize}

\newpage

% ================= PRÁCTICA 2 =================
\phantomsection
\section*{\centering PRÁCTICA 2 \\ DETERMINACIÓN DE UNA RESISTENCIA ELÉCTRICA}
\addcontentsline{toc}{section}{PRÁCTICA 2: DETERMINACIÓN DE UNA RESISTENCIA ELÉCTRICA}
 
\phantomsection
\subsection*{1. OBJETIVOS}
\addcontentsline{toc}{subsection}{1. OBJETIVOS}
\begin{itemize}
    \item Conocer los diferentes métodos experimentales que permiten determinar una resistencia eléctrica.
    \item Escoger entre estos métodos, aquel que permita determinar una resistencia de la manera más práctica y precisa.
\end{itemize}

\phantomsection
\subsection*{2. PARTE TEÓRICA}
\addcontentsline{toc}{subsection}{2. PARTE TEÓRICA}
Ver sección 2 de la primera práctica.

\phantomsection
\subsection*{3. PARTE EXPERIMENTAL}
\addcontentsline{toc}{subsection}{3. PARTE EXPERIMENTAL}
En esta parte del laboratorio, Ud. podrá determinar el valor de una resistencia eléctrica utilizando dos métodos 
experimentales, uno de medida indirecta y otro de medida directa.

\subsubsection*{3.1. MATERIAL}
\addcontentsline{toc}{subsubsection}{3.1. MATERIAL}
\begin{itemize}
    \item 1 fuente de poder de 0 – 30 v, 6 A
    \item 2 multimedidores Keithley (179, 175 ó 175-A)
    \item 1 Resistencia de carbón.
    \item 3 Resistencias variable del tipo décadas.
    \item 1 fuente estabilizada electrónicamente (Pila tipo D).
    \item Cables de conexión.
\end{itemize}

\subsubsection*{3.2. MÉTODO DE MEDIDA INDIRECTA DE UNA RESISTENCIA ELÉCTRICA}
\addcontentsline{toc}{subsubsection}{3.2. MÉTODO DE MEDIDA INDIRECTA DE UNA RESISTENCIA ELÉCTRICA}
\paragraph{3.2.1 AMPERÍMETRO Y VOLTÍMETRO.}
\addcontentsline{toc}{subsubsection}{3.2.1 AMPERÍMETRO Y VOLTÍMETRO.}
El análisis de este circuito nos indica que la corriente I que pasa por el amperímetro es la suma de la corriente 
$I_R$ que circula por el resistor R, más la corriente $I_v$ que circula por la resistencia interna del voltímetro 
$R_v$; por lo que:

\begin{equation}
    I = I_R + I_v
\end{equation}
Como $I_R = V/R$, de la misma forma que $I_v = V/R_v$ se tiene que $I = V/R + V/R_v$.
Si $R_v \gg R$, prácticamente toda la corriente pasa a través de R. Por lo tanto, $V/R_v \approx 0$ por lo que se 
puede determinar el valor de R a partir de la relación,
\begin{equation}
    R = \frac{V}{I}
\end{equation}
\begin{itemize}
    \item Arme el circuito de la figura~\ref{fig:voltimetro_fuente}.
    \item Variando la salida de la fuente, obtenga 20 valores de corriente y voltaje.
\end{itemize}
En este caso, la caída de tensión V es la suma de la caída de tensión en el amperímetro $V_A$ más la caída de tensión 
en la resistencia $V_R$. Esto es:
\begin{equation}
    V = I R_A + I R
\end{equation}
Si $R_A \ll R$, toda la caída de potencial se debe a la resistencia R, por lo tanto $I R_A \approx 0$, lo cual conduce 
nuevamente a la relación de la ecuación~\eqref{eq:ohm_basica}:
\begin{equation}
    R = \frac{V}{I}
\end{equation}
\begin{itemize}
    \item Arme el circuito de la figura~\ref{fig:voltimetro_carga}.
    \item Variando la salida de la fuente, obtenga 20 valores de corriente y voltaje.
\end{itemize}

\begin{figure}[H]
    \centering
    % --- ESQUEMA 4 ---
    \begin{minipage}[b]{0.45\textwidth}
        \centering
        \begin{circuitikz}[american, scale=0.9]
            % Fuente y Tierra
            \draw (0,0) node[ground]{} 
                  to[V, l=$E$, invert] (0,4);
            
            % Voltímetro en paralelo a la FUENTE
            \draw (0,4) -- (-2.5,4)
                  to[voltmeter, l=$V$] (-2.5,0) -- (0,0);

            % Amperímetro y Resistencia
            \draw (0,4) to[ammeter, l=$A$] (3,4)
                  to[R, l=$R$] (3,0)
                  -- (0,0);
            \node[ocirc] at (0,0) {};
        \end{circuitikz}
        \caption{Voltímetro en fuente E}
        \label{fig:voltimetro_fuente}
    \end{minipage}
    \hfill
    % --- ESQUEMA 5 ---
    \begin{minipage}[b]{0.45\textwidth}
        \centering
        \begin{circuitikz}[american, scale=0.9]
            % Fuente y Tierra
            \draw (0,0) node[ground]{} 
                  to[V, l=$E$, invert] (0,4);
            
            % Amperímetro en serie
            \draw (0,4) to[ammeter, l=$A$] (3,4)
                  to[R, l=$R$] (3,0)
                  -- (0,0);

            % Voltímetro en paralelo a la RESISTENCIA
            \draw (3,4) -- (5,4)
                  to[voltmeter, l=$V$] (5,0)
                  -- (3,0);

            \node[ocirc] at (0,0) {};
            \node[ocirc] at (3,4) {};
            \node[ocirc] at (3,0) {};
        \end{circuitikz}
        \caption{Voltímetro en carga R}\label{fig:voltimetro_carga}
    \end{minipage}
    %\caption{Circuitos para determinar la resistencia R (4 y 5)}
\end{figure}

\paragraph{3.2.2 PUENTE DE RESISTENCIAS}
\addcontentsline{toc}{subsubsection}{3.2.2 PUENTE DE RESISTENCIAS}
El puente de resistencia o puente Wheatstone, se utiliza para determinar el valor de un resistor 
desconocido por comparación con tres resistores cuyas resistencias se pueden modificar. El circuito 
de este puente, mostrado en la figura~\ref{fig:puente_wheatstone}, básicamente consta de f.e.m, 
tres resistores de precisión [$R_1, R_2$ y $R_3$], un detector V sensible a pequeñas diferencias 
de potencial y el resistor $R_x$ cuyo valor de resistencia se va a determinar.

El método consiste en balancear el puente variando las resistencias $R_1, R_2$ y $R_3$ hasta que el detector indique 
cero potencial. En estas condiciones los puntos b y d están a un mismo potencial. Por lo tanto, la caída de potencial 
a través de $R_x$ debe ser igual a la caída de potencial a través de $R_1$, esto es:
\[ I_2 R_x = I_1 R_1 \]
De la misma forma, la caída de potencial a través de $R_3$ debe ser igual a la caída de potencial a través de $R_2$, 
es decir:
\[ I_2 R_3 = I_1 R_2 \]
Si dividimos miembro a miembro ambas ecuaciones obtenemos para la resistencia desconocida el valor dado por:
\begin{equation}
    \label{eq:puente_wheatstone}
    R_x = R_3 \frac{R_1}{R_2}
\end{equation}

\begin{figure}[h!]
    \centering
    \begin{circuitikz}[american, scale=1.2, transform shape]
        % Nodos
        \node[circ] at (0,2) (A) {};
        \node[circ] at (0,-2) (B) {};
        \node[circ] at (-2,0) (C) {};
        \node[circ] at (2,0) (D) {};

        % Fuente de alimentacion
        \draw (B) -- ++(-4,0) to[V, v_=$V_s$, invert] ++(0,4) -- (A);

        % Resistencias
        \draw (A) to[R, l=$R_1$] (C);
        \draw (C) to[R, l=$R_2$] (B);
        \draw (A) to[R, l=$R_3$] (D);
        \draw (D) to[R, l=$R_x$ (desconocida)] (B);

        % Galvanometro
        \draw (C) to[ammeter, l_=G] (D);
        
        % Etiquetas de nodos
        \node[above left] at (A) {A};
        \node[below left] at (B) {B};
        \node[left] at (C) {C};
        \node[right] at (D) {D};
    \end{circuitikz}
    \caption{Diagrama de un Puente de Wheatstone.}\label{fig:puente_wheatstone}
\end{figure}

\begin{itemize}
    \item Arme el circuito mostrado en la figura~\ref{fig:puente_wheatstone}.
    \item Empezando con la sensibilidad mínima (rango máximo) del detector (voltímetro), ajuste los valores $R_1, R_2$ y $R_3$ 
    hasta obtener una indicación mínima en el detector (V = 0).
    \item Aumente la sensibilidad del detector y repita el paso anterior.
    \item Ajuste cada vez más hasta obtener la indicación cero con la máxima sensibilidad, (mínimo rango).
    \item Anote los valores de $R_1, R_2$ y $R_3$.
\end{itemize}
\textbf{RECOMENDACIÓN:} A pesar de que el ajuste a cero en el detector se logra moviendo los selectores de $R_1, R_2$ y $R_3$, 
siempre es recomendable elegir una relación $R_1/R_2$ a valores como $10^{-3}, 10^{-2}, 10^{-1}, 1, 10, 10^2$, ó $10^3$. Las 
resistencias $R_1, R_2$ y $R_3$ están constituidas por resistores variables en cuatro décadas que permiten estimar un error del 
1\% del valor ajustado. Como f.e.m, puede utilizarse una pila y como resistor $R_x$ el resistor desconocido que le suministre 
el profesor.

\subsubsection*{3.3. MÉTODO DE MEDIDA DIRECTA DE UNA RESISTENCIA ELÉCTRICA}
\addcontentsline{toc}{subsubsection}{3.3. MÉTODO DE MEDIDA DIRECTA DE UNA RESISTENCIA ELÉCTRICA}
\paragraph{3.3.1. OHMÍMETRO}
\addcontentsline{toc}{subsubsection}{3.3.1. OHMÍMETRO}
Existen multimedidores analógicos y digitales que, además de medir corriente y voltaje, están calibrados para indicar valores de 
resistencia eléctrica. Esto constituye el llamado ohmímetro.

Utilice la resistencia empleada por usted en los métodos anteriores (amperímetro y voltímetro, puente de Wheatstone) y determine el 
valor con su error, utilizando el multimedidor Keithley 175 ó 179.

\phantomsection
\subsection*{4. ACTIVIDADES}
\addcontentsline{toc}{subsection}{4. ACTIVIDADES}
\subsubsection*{4.1. LA PRECISIÓN DEL RESULTADO}
\addcontentsline{toc}{subsubsection}{4.1. LA PRECISIÓN DEL RESULTADO}
Los valores de R obtenidos mediante los métodos descritos anteriormente están afectados por un error. En las mediciones donde se 
empleó el método expuesto en la parte 3.2.1, el error depende de los errores con que son afectados los valores de V e I según las 
especificaciones de los multímetros utilizados y el rango empleado en las medidas correspondientes. Las mediciones donde se empleó 
el método desarrollado en la parte 3.2.2, partiendo del supuesto de que los cables de conexión carecen de resistencia y de que hay 
buenos contactos eléctricos en los diferentes puntos de interconexión, las fuentes de error provienen de las resistencias empleadas 
para el balanceo del puente, cuyo valor es de $\pm 1\%$ del valor ajustado. Finalmente, el error que se comete con el empleo del 
método expresado en 3.3.1, proviene del instrumento y el rango empleado, según especificaciones del manual de instrucciones.

\subsubsection*{4.2. ACTIVIDADES ADICIONALES PREVIAS Y POSTERIORES A LA REALIZACIÓN DE LA PRÁCTICA}
\addcontentsline{toc}{subsubsection}{4.2. ACTIVIDADES ADICIONALES PREVIAS Y POSTERIORES A LA REALIZACIÓN DE LA PRÁCTICA}
\begin{itemize}
    \item Revise, en la bibliografía recomendada, lo relacionado con métodos experimentales para medir resistencias eléctricas y 
    amplíe la información teórica que se ha dado aquí como introducción.
    \item Para el método 3.2.1, construya la gráfica I Vs V, aplique mínimo cuadrado y a partir de la pendiente y su error; 
    determine el valor, con su error, de la resistencia suministrada.
    \item Para el método 3.2.2, determine el valor de $R_x$ con su error, usando la ecuación~\eqref{eq:puente_wheatstone}.
    \item Para el método 3.3.1, determine el valor de $R_x$ con su error.
    \item Examine, analice y discuta, en términos de precisión, los métodos experimentales utilizados.
    \item En base a la discusión anterior y a los resultados obtenidos, concluya acerca del método que introduce menos error en la 
    determinación de una resistencia eléctrica.
\end{itemize}

\newpage

% ================= PRÁCTICA 3 =================
\phantomsection
\section*{\centering PRÁCTICA 3 \\ DETERMINACIÓN DE LA RESISTENCIA INTERNA DE UN VOLTÍMETRO, UN AMPERÍMETRO Y UNA FUENTE}
\addcontentsline{toc}{section}{PRÁCTICA 3: DETERMINACIÓN DE LA RESISTENCIA INTERNA...}
 
\phantomsection
\subsection*{1. OBJETIVO}
\addcontentsline{toc}{subsection}{1. OBJETIVO}
\begin{itemize}
    \item Determinar experimentalmente la resistencia interna de un voltímetro, un amperímetro y una fuente no 
    estabilizada electrónicamente.
\end{itemize}

\phantomsection
\subsection*{2. PARTE TEÓRICA}
\addcontentsline{toc}{subsection}{2. PARTE TEÓRICA}
Ver sección 2 de la primera práctica.

Los instrumentos de medidas, por razones intrínsecas a los mismos, tienden a ofrecer una resistencia al paso de la corriente 
eléctrica. Esta resistencia es, en general, variable de acuerdo a las posibilidades de medidas del instrumento, así como el 
tipo de corriente, (continua o variable), a la cual se conecta.

Una fuente de f.e.m., también posee una resistencia interna que determina una caída de tensión entre sus bornes. Aunque esta 
resistencia por lo general es muy pequeña, puede ser calculada mediante un montaje experimental adecuado.

\phantomsection
\subsection*{3. PARTE EXPERIMENTAL}
\addcontentsline{toc}{subsection}{3. PARTE EXPERIMENTAL}

\subsubsection*{3.1. MATERIAL}
\addcontentsline{toc}{subsubsection}{3.1. MATERIAL}
\begin{itemize}
    \item 1 multimedidor Keithley (179, 175 ó 175-A)
    \item 1 Resistencia variable del tipo décadas.
    \item 1 fuente estabilizada electrónicamente (Pila tipo D).
    \item Cables de conexión.
\end{itemize}

\subsubsection*{3.2. PROCEDIMIENTO}
\addcontentsline{toc}{subsubsection}{3.2. PROCEDIMIENTO}
\paragraph{3.2.1 RESISTENCIA INTERNA DE UN VOLTÍMETRO}
\addcontentsline{toc}{subsubsection}{3.2.1 RESISTENCIA INTERNA DE UN VOLTÍMETRO}
Puede demostrarse que la señal detectada por el voltímetro de la figura~\ref{fig:resistencia_voltimetro} es:
\begin{equation}
    V_m = \frac{R_v E}{R+R_v}
    \label{eq:res_voltimetro}
\end{equation}
Donde: R es un resistor ajustado de tal manera que no se exceda el rango del voltímetro; $R_v$ corresponde a la resistencia interna 
del voltímetro y E la salida de la fuente, la cual se determina con R=0. Así, conocidos R, E y $V_m$, se puede calcular $R_v$ para 
el rango respectivo del voltímetro.
\begin{itemize}
    \item Mida la salida de la fuente (E).
    \item Arme el circuito mostrado en la figura~\ref{fig:resistencia_voltimetro}.
    \item Seleccione el rango de 2 V en el voltímetro (179, 175 o 175-A).
    \item Ajuste un valor de R (puede ser en el orden de los M$\Omega$) y mida $V_m$.
    \item En una tabla ordene los valores de E, $V_m$ y R.
    \item Cambie el modelo del multimedidor y repita los pasos anteriores.
\end{itemize}

\begin{figure}[H]
    \centering
    \begin{circuitikz}[american]
        \draw (0,0) node[ground]{} 
              to[V, l=$E$, invert] (0,3)
              to[potentiometer, l=$R$] (3,3)
              to[voltmeter, l=$V_m$] (6,3)
              -- (6,0)
              -- (0,0);
        \node[ocirc] at (0,0) {};
    \end{circuitikz}
    \caption{Montaje Experimental para Determinar la Resistencia Interna ($R_v$) de un Voltímetro.}\label{fig:resistencia_voltimetro}
\end{figure}

\paragraph{3.2.2 RESISTENCIA INTERNA DE UN AMPERÍMETRO}
\addcontentsline{toc}{subsubsection}{3.2.2 RESISTENCIA INTERNA DE UN AMPERÍMETRO}
Puede demostrarse que la corriente, $I_m$ detectada por el amperímetro de la figura~\ref{fig:resistencia_amperimetro} es:
\begin{equation}
    I_m = \frac{E}{R+R_a}
    \label{eq:res_amperimetro}
\end{equation}
Donde: R es un resistor, ajustado de tal manera que no se exceda el rango del amperímetro; $R_a$ corresponde a la resistencia interna 
del medidor y E es la salida de la fuente.

Así conociendo E, R, e $I_m$, se puede determinar el valor de $R_a$, para el rango respectivo del amperímetro.
\begin{itemize}
    \item Mida la salida de la fuente (E).
    \item Arme el circuito mostrado en la figura~\ref{fig:resistencia_amperimetro}.
    \item Seleccione el rango de 200 $\mu$A en el amperímetro.
    \item Ajuste un valor de R, tal que no exceda el rango del amperímetro y mida $I_m$.
    \item En una tabla ordene los valores de E, $I_m$ y R.
    \item Cambie el rango del multimedidor (2 mA) y repita los pasos anteriores.
\end{itemize}

\begin{figure}[H]
    \centering
    \begin{circuitikz}[american]
        \draw (0,0) node[ground]{} 
              to[V, l=$E$, invert] (0,3)
              to[potentiometer, l=$R$] (3,3)
              to[ammeter, l=$A_m$] (6,3)
              -- (6,0)
              -- (0,0);
        \node[ocirc] at (0,0) {};
    \end{circuitikz}
    \caption{Montaje Experimental para Determinar la Resistencia Interna ($R_a$) de un Amperímetro.}\label{fig:resistencia_amperimetro}
\end{figure}

\paragraph{3.2.3 DETERMINACIÓN DE LA RESISTENCIA INTERNA DE UNA FUENTE}
\addcontentsline{toc}{subsubsection}{3.2.3 DETERMINACIÓN DE LA RESISTENCIA INTERNA DE UNA FUENTE}
El voltaje $V_m$ registrado por el voltímetro de la figura~\ref{fig:resistencia_fuente}, está dado por la ecuación:
\begin{equation}
    V_m = \frac{E R_e}{R_s + R_e}
    \label{eq:res_fuente_vm}
\end{equation}
Donde: $R_e$ representa la resistencia equivalente, para el conjunto formado por R y $R_v$ 
(resistencia interna el voltímetro), las cuales se encuentran conectadas en paralelo. Por tanto,
\begin{equation}
    R_e = \frac{R R_v}{R+R_v}
    \label{eq:res_fuente_re}
\end{equation}
La ecuación~\eqref{eq:res_fuente_vm} puede reescribirse para observar la relación lineal que hay entre 
$R_e/V_m$ y $R_e$; o sea,
\begin{equation}
    \frac{R_e}{V_m} = \frac{1}{E} R_e + \frac{R_s}{E}
    \label{eq:res_fuente_lineal}
\end{equation}
\begin{itemize}
    \item Mida la salida de la fuente (E).
    \item Arme el circuito mostrado en la figura~\ref{fig:circuito_diodo}.
    \item Variando la resistencia obtenga 20 valores de $V_m$ y R (Ordénelos en una tabla).
\end{itemize}

\begin{figure}[H]
    \centering
    \begin{circuitikz}[american]
        \draw (0,0) node[ground]{} to[V, l=$E$, invert] (0,3)
              to[R, l=$R_s$] (3,3)
              to[potentiometer, l=$R$] (3,0) -- (0,0);
        \draw (3,3) -- (5,3) to[voltmeter, l=$V$] (5,0) -- (3,0);
        \node[ocirc] at (0,0) {};
        \node[ocirc] at (3,3) {};
        \node[ocirc] at (3,0) {};
    \end{circuitikz}
    
    \caption{Montaje Experimental para Determinar la Resistencia Interna $R_s$, de una Fuente.}
    \label{fig:resistencia_fuente}
\end{figure}

\phantomsection
\subsection*{4. ACTIVIDADES}
\addcontentsline{toc}{subsection}{4. ACTIVIDADES}
\subsubsection*{4.1. RESISTENCIA INTERNA DE UN VOLTÍMETRO}
\addcontentsline{toc}{subsubsection}{4.1. RESISTENCIA INTERNA DE UN VOLTÍMETRO}
\begin{itemize}
    \item Deduzca la expresión~\eqref{eq:res_voltimetro}.
    \item Determine, con su error, la resistencia interna de los voltímetros digitales 179 y 175 en el rango de 2V.
    \item Determine la desviación porcentual del valor medido por Ud., $R_v$, respecto al valor indicado por el fabricante 
    $R_F$ (10M$\Omega$ para 179 y 11M$\Omega$ para 175), mediante la expresión:
    \begin{equation}
        E\% = \frac{|R_F - R_v|}{R_F} \times 100
        \label{eq:error_porcentual}
    \end{equation}
\end{itemize}

\subsubsection*{4.2. RESISTENCIA INTERNA DE UN AMPERÍMETRO}
\addcontentsline{toc}{subsubsection}{4.2. RESISTENCIA INTERNA DE UN AMPERÍMETRO}
\begin{itemize}
    \item Deduzca la expresión~\eqref{eq:res_amperimetro}.
    \item Determine el valor con su error, de la resistencia interna $R_a$ de los multimedidores Keithley 179 ó 175 en 
    los rangos de 200$\mu$A y 2mA.
    \item Determine la desviación porcentual del valor medido por Ud. $R_a$, con respecto al valor indicado por el 
    fabricante $R_f$ (1k$\Omega$ para 200$\mu$A y 100$\Omega$ para 2 mA), mediante la ecuación~\eqref{eq:error_porcentual}.
\end{itemize}

\subsubsection*{4.3. RESISTENCIA INTERNA DE UNA FUENTE}
\addcontentsline{toc}{subsubsection}{4.3. RESISTENCIA INTERNA DE UNA FUENTE}
\begin{itemize}
    \item Deduzca las ecuaciones~\eqref{eq:res_fuente_vm}, \eqref{eq:res_fuente_re} y~\eqref{eq:res_fuente_lineal}.
    \item Para cada valor de R, determine $R_e$, usando la ecuación~\eqref{eq:res_fuente_re}.
    \item Grafique $R_e/V_m$ Vs $R_e$. Aplique mínimo cuadrado y determine la pendiente y el corte con sus respectivos errores.
    \item A partir del valor del corte, determine el valor con su error, de la resistencia interna de la fuente suministrada.
\end{itemize}

\newpage

% ================= PRÁCTICA 4 =================
\phantomsection
\section*{\centering PRÁCTICA 4 \\ CIRCUITOS ELÉCTRICOS}
\addcontentsline{toc}{section}{PRÁCTICA 4: CIRCUITOS ELÉCTRICOS}

\phantomsection
\subsection*{1. OBJETIVOS}
\addcontentsline{toc}{subsection}{1. OBJETIVOS}
\begin{itemize}
    \item Medir corrientes y voltajes en circuitos conectados en serie, paralelo y mixto.
    \item Aplicar las leyes de Kirchhoff en redes eléctricas.
\end{itemize}

\phantomsection
\subsection*{2. METODOLOGÍA}
\addcontentsline{toc}{subsection}{2. METODOLOGÍA}
Se usa una red en la cual hay conectadas varias baterías y cajas de resistencias. Se determinan las intensidades de corriente 
utilizando valores conocidos de las resistencias y de los voltajes de las baterías.

\phantomsection
\subsection*{3. PARTE TEÓRICA}
\addcontentsline{toc}{subsection}{3. PARTE TEÓRICA}
Dos o más resistencias conectadas de modo que la misma carga fluya a través de cada una de ellas se dice que están conectadas en 
serie; por ella debe circular la misma corriente i.

\begin{figure}[H]
    \centering
    \begin{circuitikz}[american]
        \draw (0,0) node[circ, label=below:a] {} to[R, l=$R_1$, v=$V_1$, i=$I_1$] (3,0) to[R, l=$R_2$, v=$V_2$, i=$I_2$] 
        (6,0) node[circ, label=below:b] {};
    \end{circuitikz}
    \caption{Dos resistencias conectadas en serie por las cuales circula una corriente de intensidad I.}
    \label{fig:serie_res}
\end{figure}

Para este tipo de conexión deben cumplirse las siguientes ecuaciones:
\begin{subequations}
\label{eq:serie}
\begin{align}
    \text{Corriente: } & I = I_1 = I_2 \label{eq:serie_i} \\
    \text{Voltaje: }   & V_{ab} = V_1 + V_2 \label{eq:serie_v} \\
    \text{Resistencia:} & R_{eq} = R_1 + R_2 \label{eq:serie_r}
\end{align}
\end{subequations}

Dos resistencias unidas como en la figura~\ref{fig:paralelo_res} se dice que están conectadas en paralelo. Sea I la 
corriente que va desde 'a' hasta 'b'. 
En el punto 'a', la corriente se divide en dos, $I_1$ que circula por $R_1$ e $I_2$ por $R_2$.

\begin{figure}[H]
    \centering
    \begin{circuitikz}[american, scale=1.2] 
        % 1. Dibujar la fuente de corriente I_s en el lado izquierdo
        % Usamos (0, 0) como centro, el diagrama se verá mejor centrado verticalmente
        
        % Fuente de corriente independiente (isource)
        \draw (-2, 0) to[isource, l_=$I$, *-*] (0, 0);
        
        \draw (0, 1) -- (0, -1); % Línea vertical de la izquierda

        % Etiqueta del nodo 'a' (el punto donde la corriente entra al paralelo)
        \draw (0, 0) node[circ, label=right:a]{};
        
        % 2. Corriente que Circula por Rama superior (R1)
        \draw (0, 1)
              to[R, l_=$R_1$, i>_=$I_1$] (4, 1);
        
        % 3. Corriente que Circula por Rama inferior (R2)
        \draw (0, -1)
              to[R, l_=$R_2$, i>_=$I_2$] (4, -1);

        % 4. Conectar el lado derecho y etiquetar el nodo b
        % Línea vertical de la derecha
        \draw (4, 1) -- (4, -1); 
        
        % Etiqueta del nodo 'b' (el punto donde la corriente sale del paralelo)
        \draw (4, 0) node[circ, label=left:b]{};
        
        % 5. Cables de salida de la corriente total (I_s) por el nodo b (opcional, por claridad)
        \draw (4, 0) -- (5, 0);
        %\draw (4, -1) -- (5, -1);
        
        \draw[->] (5, 0) -- (5.1, 0) node[right, pos=1] {$I$};
    
    \end{circuitikz}
    \caption{Circuito de dos resistencias en paralelo $R_1$ y $R_2$ alimentado por una fuente de corriente $I$. 
    La corriente $I$ entra por el nodo 'a' y se divide en $I_1$ e $I_2$, saliendo por el nodo 'b'.}
    \label{fig:paralelo_res}
\end{figure}

Para este tipo de conexión deben cumplirse las siguientes ecuaciones:
\begin{subequations}
\label{eq:paralelo}
\begin{align}
    \text{Corriente: } & I = I_1 + I_2 \label{eq:paralelo_i} \\
    \text{Voltaje: }   & V_{ab} = V_1 = V_2 \label{eq:paralelo_v} \\
    \text{Resistencia:} & R_{eq} = \frac{R_1 R_2}{R_1 + R_2} \label{eq:paralelo_r}
\end{align}
\end{subequations}

A veces no es posible reducir muchas de las redes importantes de electrónica a simples combinaciones serie-paralelo, 
por lo que deben utilizarse métodos analíticos más avanzados. A tal fin, resultan muy útiles las conocidas leyes de 
Kirchhoff. Estas leyes pueden ser enunciadas de la siguiente manera:
\begin{description}
    \item[Primera ley:] En cualquier nodo de un circuito, la suma de las corrientes que llegan a él es igual a la 
    suma de las corrientes que salen del mismo, es decir, la suma algebraica de las corrientes en un nodo es igual 
    a cero. ecuación~\eqref{eq:kcl}.

    \begin{equation}
        \sum I = 0
        \label{eq:kcl}
    \end{equation}

    \item[Segunda ley:] En una malla, la suma algebraica de los productos IR y de las f.e.m que se encuentran en un 
    recorrido es igual a cero. ecuación~\eqref{eq:kvl}.

    \begin{equation}
        \sum IR + \sum \epsilon_i = 0
        \label{eq:kvl}
\end{equation}

\end{description}

\subsubsection*{Reglas para la Aplicación Correcta de las Leyes de Kirchhoff:}
\addcontentsline{toc}{subsubsection}{Reglas para la Aplicación Correcta de las Leyes de Kirchhoff:}
\begin{itemize}
    \item \textbf{Nodo:} punto donde se unen tres o más conductores.
    \item \textbf{Malla:} Cualquier recorrido seguido de un circuito cerrado.
    \item \textbf{Signos algebraicos e intensidades:} se acostumbra tomar como positivas las corrientes que llegan a 
    un nodo y negativas las que salen de él.
    \item \textbf{Signos algebraicos de voltajes:} Al recorrer una malla en cualquier dirección se encontrarán resistencias 
    y baterías. En las resistencias las diferencias de potencial son positivas si se atraviesan de los puntos de mayor a menor 
    potencial y negativas si se atraviesan de menor a mayor potencial.
    \item \textbf{La fem, batería o pila ($\epsilon$):} es positiva si se recorre del borne (-) al borne (+) y negativa si se 
    recorre de (+) a (-).
    \item Dibuje un diagrama de la red, claro y de tamaño suficiente. Use letras y símbolos para cada corriente, resistencia o 
    voltaje desconocido.
    \item Indique el sentido de cada corriente con una flecha. Si el sentido no lo conoce, suponga momentáneamente un sentido en 
    cualquier dirección. Si supone errado el sentido, la corriente le aparecerá con signo (-) al resolver las ecuaciones.
    \item Coloque un signo (+) al terminal positivo de cada batería y al extremo de cada resistencia por donde usted supone que 
    la corriente entra.
    \item Si tiene 'n' nodos entonces habrá (n-1) ecuaciones independientes.
    \item Si 'm' es el número de mallas, entonces usted tendrá (m-1) ecuaciones independientes.
\end{itemize}

\phantomsection
\subsection*{4. PARTE EXPERIMENTAL}
\addcontentsline{toc}{subsection}{4. PARTE EXPERIMENTAL}
\subsubsection*{4.1. MATERIAL}
\addcontentsline{toc}{subsubsection}{4.1. MATERIAL}
\begin{itemize}
    \item 1 Tablero para redes.
    \item Batería de 10 volt.
    \item 3 Resistencias de carbón.
    \item 1 multimedidor Keithley modelo 179 o 175.
\end{itemize}

\subsubsection*{4.2. PROCEDIMIENTO}
\addcontentsline{toc}{subsubsection}{4.2. PROCEDIMIENTO}
\paragraph{4.2.1. RESISTENCIAS EN SERIE}
\addcontentsline{toc}{subsubsection}{4.2.1. RESISTENCIAS EN SERIE}
\begin{itemize}
    \item Mida con el voltímetro la fem $\epsilon$ y con el Ohmímetro los valores de $R_1$ y $R_2$. Organice estos 
    datos en una tabla.
    \item Arme el circuito mostrado en la figura~\ref{fig:exp_serie}.
    \item Mida las corrientes ($I, I_1$ e $I_2$) con el amperímetro.
    \item Mida la diferencia de potencial entre los extremos de cada resistencia.
\end{itemize}

\begin{figure}[H]
    \centering
    \begin{circuitikz}[american]
        \draw (0,0) 
            to[V, l=$\epsilon$, invert] (0,3) 
            to[R, l=$R_1$, i=$I_1$] (3,3) 
            to[R, l=$R_2$, i=$I_2$] (6,3) -- (6,0) -- (0,0);
        \draw (0,2) to[short, i=I] (0,3);
        \node[ocirc] at (0,0) {};
        \draw (0,0) node[ground]{};
    \end{circuitikz}
    \caption{Circuito Experimental: fem ($\epsilon$) en Serie con $R_1$ y $R_2$.}
    \label{fig:exp_serie}
\end{figure}

\paragraph{4.2.2. RESISTENCIAS EN PARALELO}
\addcontentsline{toc}{subsubsection}{4.2.2. RESISTENCIAS EN PARALELO}
\begin{itemize}
    \item Mida con el voltímetro la f.e.m. $\epsilon$ y con el Ohmímetro los valores de $R_1$ y $R_2$. Organice estos 
    datos en una tabla.
    \item Arme el circuito mostrado en la figura~\ref{fig:exp_paralelo}.
    \item Mida la corriente total (I) la que circula por cada resistencia ($I_1$ e $I_2$).
    \item Mida la diferencia de potencial entre los extremos de cada resistencia.
\end{itemize}

\begin{figure}[H]
    \centering
    \begin{circuitikz}[american]
        % Rama 1 (Izquierda): Fuente de fem
        % Dibujamos la fuente de (0,0) a (0,4)
        \draw (0,0) to[V, l=$\epsilon$, invert] (0,4);
        
        % Línea superior (Rail positivo)
        % Conecta la fuente con las ramas de las resistencias
        \draw (0,4) -- (4,4);
        
        % Rama 2 (Centro): Resistencia R1
        % Conectada entre el rail superior e inferior
        % El nodo[circ] dibuja el punto de conexión (nodo)
        \draw (2,4) node[circ]{} 
              to[R, l=$R_1$, i=$I_1$] (2,0) 
              node[circ]{};
        
        % Rama 3 (Derecha): Resistencia R2
        \draw (4,4) to[R, l=$R_2$, i=$I_2$] (4,0);
        
        % Línea inferior (Rail negativo / Tierra)
        % Cierra el circuito volviendo a la fuente
        \draw (4,0) -- (0,0);

        \draw (0,3) to[short, i=I] (0,4);
        
        % Opcional: Agregar tierra en el negativo de la fuente
        \draw (0,0) node[ground]{};
    \end{circuitikz}
    \caption{Circuito Experimental, fem ($\epsilon$) en Paralelo con $R_1$ y $R_2$.}
    \label{fig:exp_paralelo}
\end{figure}

\paragraph{4.2.3. RED ELÉCTRICA}
\addcontentsline{toc}{subsubsection}{4.2.3. RED ELÉCTRICA}
\begin{itemize}
    \item Disponga las resistencias y pilas en el tablero tal como se muestra en la figura~\ref{fig:exp_red}.
    \item Mida con su error los valores de $\epsilon_1, \epsilon_2, \epsilon_3, R_1, R_2$ y $R_3$.
    \item Arme el correspondiente circuito.
    \item Mida la corriente de cada rama ($I_1, I_2,$ e $I_3$); para esto, rompa una rama del circuito e intercale el 
    amperímetro (en serie). Haga revisar las conexiones por el profesor.
    \item Mida para cada malla, la diferencia de potencial entre los extremos de cada resistencia.
\end{itemize}

\begin{figure}[H]
    \centering
    \begin{circuitikz}[american]
        
        % Malla 1 (Izquierda)
        \draw (0,0) 
            to[V, l=$\epsilon_1$, invert] (0,2) -- (0,4);
        
        \draw (0,4)
            to[R, l=$R_1$, i=$I_1$] (4,4); % R1 hacia el nodo central superior

        \draw (4,4) 
            to[V, l=$\epsilon_2$] (4, 1)
            to[R, l=$R_2$, i<=$I_2$] (4,-2); % R2: En serie, conectada al nodo inferior

        % Malla 2 (Derecha)
        \draw (4,4) 
            to[R, l=$R_3$, i<=$I_3$] (8,4) % R3 desde el nodo central superior
            to[V, l=$\epsilon_3$] (8,-2); % E3 (polaridad estándar)

        % Cierre de la Malla 2 (Rail Inferior Común)
        \draw (8,-2) -- (4,-2) -- (4,-2) -- (0,-2) -- (0,0); 

        % Añadir nodos explícitos en las uniones para mayor claridad
        %\node[ocirc] at (4,3) {};
        \node[ocirc] at (4,-2) {}; % Punto de conexión de R2 al rail inferior
        \node[ocirc] at (4,4) {};

        % Opcional: Agregar tierra en el negativo de la fuente
        \draw (4,-2) node[ground]{};
    \end{circuitikz}
    \caption{Red Eléctrica}
    \label{fig:exp_red}
\end{figure}

\phantomsection
\subsection*{5. ACTIVIDADES}
\addcontentsline{toc}{subsection}{5. ACTIVIDADES}
\subsubsection*{5.1. RESISTENCIAS EN SERIE}
\addcontentsline{toc}{subsubsection}{5.1. RESISTENCIAS EN SERIE}
\begin{itemize}
    \item Calcule la corriente que circula por el circuito y compárela con el valor medido.
    \item Verifique si se cumple la ecuación~\eqref{eq:serie_v}.
\end{itemize}

\subsubsection*{5.2. RESISTENCIAS EN PARALELO}
\addcontentsline{toc}{subsubsection}{5.2. RESISTENCIAS EN PARALELO}
\begin{itemize}
    \item Calcule la corriente total que circula por el circuito y las corrientes $I_1$ e $I_2$. Compárelas con el valor medido.
    \item Verifique si se cumple la ecuación~\eqref{eq:paralelo_i}.
    \item Verifique si se cumple la ecuación~\eqref{eq:paralelo_v}.
\end{itemize}

\subsubsection*{5.3. RED ELÉCTRICA}
\addcontentsline{toc}{subsubsection}{5.3. RED ELÉCTRICA}
\begin{itemize}
    \item Aplique las leyes de Kirchhoff. Resuelva las ecuaciones en las incógnitas $I_1, I_2,$ e $I_3$. 
    (Debe traerlo resuelto antes de entrar al laboratorio).
    \item Con los valores medidos de $\epsilon_1, \epsilon_2, \epsilon_3, R_1, R_2$ y $R_3$. Calcule los 
    valores de $I_1, I_2,$ e $I_3$.
    \item Calcule el porcentaje de error entre los valores experimentales obtenidos para las corrientes y 
    los calculados por las leyes de Kirchhoff.
    \item Verifique si se cumple la ecuación~\eqref{eq:kcl}.
    \item Verifique si se cumple la ecuación~\eqref{eq:kvl}.
\end{itemize}

\end{document}